\sect{13. Catálogo de conjuros y oraciones}
El libro identifica los hechizos de mago y el símbolo sagrado, los conjuros divinos: son los dibujos de sus herramientas en el juego, no iconos exclusivos de cada conjuro.

Las tablas distinguen el \textbf{nivel del conjuro} del nivel del personaje. Para preparar una entrada necesitas acceso a ese nivel, un hueco disponible y, si eres mago, conocerla en tu libro. El paladín dispone de una selección propia de oraciones; el menú muestra las que puede preparar.

\textbf{Cómo leer los dados.} 1d8 significa de 1 a 8; 2d8 + 1 significa tirar dos dados de ocho caras y sumar uno (3--17). Son valores base: las defensas, resistencias y tiradas de salvación pueden cambiar el resultado. Una curación nunca supera los PV máximos. Los efectos de protección repetidos no deben darse por acumulativos.

Los personajes multiclase usan el nivel de la clase que lanza el conjuro.

\refspace
\subsection*{Magia de mago}

\refspace
\subsubsection*{Nivel 1}

\begin{reference}{Nombre}{Efecto, uso y límites}
\itemname{spellbook}{Armadura} & Sobre un personaje. Ofrece una CA base de 6, ajustada por destreza, si mejora su protección. Termina cuando el personaje acumula daños por un total de 8 + nivel del lanzador; no absorbe ese daño.\\[4pt]
\itemname{spellbook}{Manos ardientes} & Ataque cercano contra los monstruos de la casilla de delante: 1d3 + 2 por nivel del lanzador. Sitúate frente al enemigo.\\[4pt]
\itemname{spellbook}{Detectar magia} & Mientras dura, los iconos de objetos mágicos se tiñen de azul en el inventario, las manos y el cursor. Al cogerlos también aparece [MAGICO] junto al nombre. No los identifica; Leer magia revela sus propiedades.\\[4pt]
\itemname{spellbook}{Dardo mágico} & Proyectil frontal. Su daño parte de 1d4 + 1 y aumenta con el nivel mediante impactos adicionales. Es práctico desde la retaguardia; las paredes interrumpen la trayectoria.\\[4pt]
\itemname{spellbook}{Leer magia} & Identifica, sin consumirlo, el objeto que llevas en el cursor. Después de lanzarlo se muestra su nombre completo; funciona también con pociones, varitas y anillos. Está disponible para preparar con cualquier mago.\\[4pt]
\itemname{spellbook}{Escudo} & Protección del lanzador. Limita su CA a 4 si era peor y detiene los dardos mágicos que le alcanzan. No es inmunidad a todos los ataques ni a toda la magia.\\[4pt]
\itemname{spellbook}{Golpe eléctrico} & Necesita una mano libre. Crea un ataque mágico en ella, de 1d8 + nivel. Después debes usar esa mano contra un enemigo cercano; no sale un proyectil al elegir el conjuro.\\[4pt]
\end{reference}

\refspace
\subsubsection*{Nivel 2}

\begin{reference}{Nombre}{Efecto, uso y límites}
\itemname{spellbook}{Invisibilidad} & Sobre un personaje. Dificulta que los monstruos lo alcancen, pero no lo hace invulnerable. Usar una mano o lanzar magia ofensiva puede romper la invisibilidad.\\[4pt]
\itemname{spellbook}{Desbloquear} & Abre la puerta ordinaria que tienes delante y libera las puertas que podrían forzarse físicamente. No altera paredes ni cerraduras mágicas o mecanismos de puzle. Está disponible para preparar con cualquier mago que tenga huecos de nivel 2.\\[4pt]
\itemname{spellbook}{Flecha ácida de Melf} & Proyectil frontal: 2d4 por cada tres niveles completos del lanzador, con un mínimo de una tirada. Aplica el daño al impactar, sin daño continuo en turnos posteriores.\\[4pt]
\itemname{spellbook}{Nube ponzoñosa} & Crea una nube en la casilla abierta de delante. Puede inmovilizar brevemente a monstruos y a miembros del grupo que entren en ella. Los no muertos son inmunes. Espera a que se disipe antes de avanzar.\\[4pt]
\end{reference}

\refspace
\subsubsection*{Nivel 3}

\begin{reference}{Nombre}{Efecto, uso y límites}
\itemname{spellbook}{Disolver magia} & Termina los efectos mágicos de un personaje, incluida la invisibilidad y las protecciones beneficiosas. No sustituye a Neutralizar veneno ni a Quitar parálisis.\\[4pt]
\itemname{spellbook}{Bola de fuego} & Proyectil explosivo: 1d6 por nivel del lanzador, antes de las defensas del objetivo. Comprueba la dirección y deja espacio para maniobrar.\\[4pt]
\itemname{spellbook}{Flecha ardiente} & Proyectil frontal de 3d10; pasa a 6d10 al alcanzar el nivel 10 del lanzador. No requiere una flecha física ni un arco.\\[4pt]
\itemname{spellbook}{Velocidad} & Reduce la recuperación total de las armas de unos 3 a 1,5 segundos para los miembros conscientes del grupo durante aproximadamente 30 segundos por (nivel del mago + 3). Si todos los miembros conscientes están acelerados, también se reduce a la mitad la espera de repetición al mantener pulsada una tecla de movimiento.\\[4pt]
\itemname{spellbook}{Paralizar / Bloquear persona} & Proyectil que intenta detener enemigos de tipo humanoide. La resistencia del objetivo puede impedirlo; no sirve contra cualquier criatura.\\[4pt]
\itemname{spellbook}{Invisibilidad 3 m} & Aplica invisibilidad al grupo vivo. Reduce la exposición a los golpes, pero el grupo puede seguir recibiendo daño. Las acciones ofensivas pueden revelar a quien las realiza.\\[4pt]
\itemname{spellbook}{Rayo eléctrico} & Ataque en línea recta, de 1d6 por nivel antes de las defensas. Alinea el pasillo y la orientación del grupo con los enemigos.\\[4pt]
\itemname{spellbook}{Toque vampírico} & Necesita una mano libre y crea un ataque de contacto, con 1d6 por cada dos niveles. Usa después esa mano. Cuando el golpe causa daño, devuelve esa cantidad de PV al atacante; el ataque se gasta al acertar.\\[4pt]
\end{reference}

\refspace
\subsubsection*{Nivel 4}

\begin{reference}{Nombre}{Efecto, uso y límites}
\itemname{spellbook}{Terror} & Intenta hacer huir a los monstruos de la casilla de delante. Puede fallar contra criaturas resistentes. Úsalo para abrir espacio para retirarte.\\[4pt]
\itemname{spellbook}{Tormenta de hielo} & Al alcanzar un objetivo, inflige 3d10 en su casilla y las cuatro adyacentes, formando una cruz. También puede herir al grupo si queda dentro del área. El daño base no depende del nivel del lanzador.\\[4pt]
\itemname{spellbook}{Piel de roca} & Protege a un personaje frente a 1d4 + la mitad del nivel del mago (redondeada hacia abajo) golpes físicos cuerpo a cuerpo que acierten. Cada golpe protegido gasta una carga. No detiene daño de conjuros y puede disiparse.\\[4pt]
\end{reference}

\refspace
\subsubsection*{Nivel 5}

\begin{reference}{Nombre}{Efecto, uso y límites}
\itemname{spellbook}{Nube ácida} & Crea una nube venenosa en la casilla abierta de delante. Puede matar a los enemigos más débiles y herir a los más fuertes. Los no muertos son inmunes. La nube también puede herir a los miembros del grupo que entren en ella.\\[4pt]
\itemname{spellbook}{Cono helado} & Ataca un abanico de casillas frente al grupo. Daño base: 1d4 + 1 por nivel; los enemigos pueden reducirlo con sus defensas. Es útil contra varios adversarios.\\[4pt]
\itemname{spellbook}{Paralizar monstruo} & Versión de inmovilización dirigida a monstruos, no solo humanoides. El proyectil debe llegar al objetivo y su resistencia puede impedir el efecto.\\[4pt]
\end{reference}

\refspace
\subsection*{Magia divina: clérigos y paladines}

\refspace
\subsubsection*{Nivel 1}

\begin{reference}{Nombre}{Efecto, uso y límites}
\itemname{holy-symbol}{Bendición} & Mejora en 1 la tirada para acertar del grupo durante su efecto. No se acumula lanzándola repetidamente sobre un grupo ya bendecido.\\[4pt]
\itemname{holy-symbol}{Curación leve} & Cura 1d8 puntos de vida al personaje elegido, sin superar su máximo. No elimina veneno ni resucita a un muerto.\\[4pt]
\itemname{holy-symbol}{Herida leve} & Ataque de contacto de 1d8. El lanzador tiene que estar en la primera fila y debe haber un enemigo delante; el intento puede fallar.\\[4pt]
\itemname{holy-symbol}{Detectar magia} & Mientras dura, los iconos de objetos mágicos se tiñen de azul en el inventario, las manos y el cursor. Al cogerlos también aparece [MAGICO] junto al nombre. No los identifica; Leer magia revela sus propiedades.\\[4pt]
\itemname{holy-symbol}{Protección contra el mal} & Protege a un personaje: los ataques de monstruos contra él sufren una penalización de 2. Sigue siendo necesario curarlo y evitar quedar rodeado.\\[4pt]
\itemname{holy-symbol}{Imposición de manos} & Facultad del paladín: cura 2 puntos de vida por nivel del lanzador. Se recupera con el descanso; no es una oración normal disponible para todos los clérigos.\\[4pt]
\end{reference}

\refspace
\subsubsection*{Nivel 2}

\begin{reference}{Nombre}{Efecto, uso y límites}
\itemname{holy-symbol}{Ayuda} & Concede 1d8 PV temporales que absorben el daño primero. Los puntos que sobren desaparecen al terminar el conjuro. Volver a lanzar Ayuda mientras está activa no suma más puntos.\\[4pt]
\itemname{holy-symbol}{Hoja llameante} & Necesita una mano libre. Crea un arma temporal de 1d4 + 4, que se utiliza como un arma cuerpo a cuerpo.\\[4pt]
\itemname{holy-symbol}{Paralizar / Bloquear persona} & Intenta inmovilizar humanoides mediante un proyectil. La resistencia puede anularlo. Su uso es equivalente al de la versión de mago.\\[4pt]
\itemname{holy-symbol}{Retrasar veneno} & Sobre un personaje envenenado. Suspende temporalmente el daño del veneno, pero no lo cura. Busca un antídoto antes de que termine.\\[4pt]
\end{reference}

\refspace
\subsubsection*{Nivel 3}

\begin{reference}{Nombre}{Efecto, uso y límites}
\itemname{holy-symbol}{Crear comida} & Llena la reserva de comida de todos los personajes vivos. No restaura puntos de vida ni elimina enfermedades o veneno.\\[4pt]
\itemname{holy-symbol}{Disipar magia} & Termina los efectos mágicos de un personaje, incluida la invisibilidad y las protecciones beneficiosas. No sustituye a Neutralizar veneno ni a Quitar parálisis.\\[4pt]
\itemname{holy-symbol}{Atuendo mágico} & Mejora la clase de armadura a 5, o a un valor menor si quien lo lanza tiene más experiencia. Conserva una protección superior que el personaje ya tuviera.\\[4pt]
\itemname{holy-symbol}{Oración} & Da +1 al ataque del grupo y penaliza en 1 los ataques enemigos. El +1 al ataque no se suma al de Bendición. No cura ni detiene el tiempo.\\[4pt]
\itemname{holy-symbol}{Quitar parálisis} & Libera hasta cuatro personajes paralizados del grupo. No es un antídoto ni una resurrección.\\[4pt]
\end{reference}

\refspace
\subsubsection*{Nivel 4}

\begin{reference}{Nombre}{Efecto, uso y límites}
\itemname{holy-symbol}{Curación grave} & Cura 2d8 + 1 puntos de vida al personaje elegido, hasta su máximo. Reserva las curas grandes para heridas importantes.\\[4pt]
\itemname{holy-symbol}{Herida grave} & Ataque de contacto de 2d8 + 1. Requiere un lanzador en primera fila y un enemigo al alcance.\\[4pt]
\itemname{holy-symbol}{Neutralizar veneno} & Elimina el envenenamiento del personaje elegido. No repone los puntos de vida que ya haya perdido; cúralos por separado.\\[4pt]
\itemname{holy-symbol}{Protección contra el mal 3 m} & Extiende la protección contra ataques de monstruos a los personajes vivos del grupo. Úsala antes de entrar en una zona peligrosa.\\[4pt]
\itemname{holy-symbol}{Protección eléctrica} & No se puede preparar.\\[4pt]
\end{reference}

\refspace
\subsubsection*{Nivel 5}

\begin{reference}{Nombre}{Efecto, uso y límites}
\itemname{holy-symbol}{Curación crítica} & Cura 3d8 + 3 puntos de vida al personaje elegido, sin superar su máximo. No revive a un muerto.\\[4pt]
\itemname{holy-symbol}{Herida crítica} & Ataque de contacto de 3d8 + 3. El lanzador debe estar en la primera fila y necesita un enemigo delante.\\[4pt]
\itemname{holy-symbol}{Columna de fuego} & Ataque mágico de 6d8 al impactar, antes de las defensas del enemigo. Es una opción ofensiva de nivel alto.\\[4pt]
\itemname{holy-symbol}{Resucitar} & Devuelve a un personaje muerto a 1 punto de vida. Los elfos son inmunes. Cura al resucitado antes de volver a combatir.\\[4pt]
\end{reference}
